\documentclass{standalone}
\usepackage{tikz}
\usepackage{amsmath}
\usetikzlibrary{patterns,decorations.markings,backgrounds,arrows.meta}

\usepackage[T1]{fontenc}
\usepackage{newtxtext}
\usepackage{newtxmath}

% jfm.cls maps the upper-case Greek letters onto their slanted (var) forms,
% so \Omega is italic in the paper; do the same here.
\let\Omega\varOmega

% ==================================================================
%  Three-dimensional view of the computational domain.
%  Oblique (cabinet-type) projection: x and y are drawn undistorted,
%  so the plane z = 0 is exactly the two-dimensional slice, and the
%  spanwise axis z is drawn along (zx,zy).  The domain is periodic
%  in z, with period b.
% ==================================================================

% ------------------------------------------------------------------
%  Geometry (streamwise/wall-normal lengths anisotropically rescaled)
% ------------------------------------------------------------------
\def\xin{-20}   % inlet
\def\xo{55}     % outlet
\def\d{4}       % gap depth
\def\w{15}      % gap width
\def\h{20}      % domain height
\def\scale{0.3}

\def\aa{\xin*\scale}
\def\dd{\d*\scale}
\def\ww{\w*\scale}
\def\oo{\xo*\scale}
\def\hh{\h*\scale}

% drawn spanwise extent (b = 4d in the simulations; strongly
% exaggerated here with respect to the streamwise lengths)
\def\zb{6.2}

% oblique projection of the spanwise axis
\def\zx{0.46}
\def\zy{0.33}
\pgfmathsetmacro{\zangle}{atan(\zy/\zx)}   % its angle on the page

% Final width = 13.5 cm
\def\geomscale{0.5089}
\pgfmathsetmacro{\zxcm}{\zx*\geomscale}
\pgfmathsetmacro{\zycm}{\zy*\geomscale}

\newcommand{\fsize}{\fontsize{8.7}{10.44}\selectfont}
\newcommand{\fsizebig}{\fontsize{12}{14.4}\selectfont}

\def\surfcol{black!6}    % wall surface (lit from above)
\def\floorcol{black!13}  % gap floor
\def\facecol{black!28}   % upstream face inside the gap
\def\edgecol{black!55}   % edges of the receding surfaces
\def\ghostcol{black!45}  % inlet profile at z = b

% Blasius profile at a given spanwise station:
%   #1 = z,  #2 = style of the profile,  #3 = style of the arrows
\newcommand{\blprofile}[3]{%
	\draw[#2]
	plot[smooth,tension=0.5] coordinates {
			(-2,0,#1)
			(-1.8,0.075,#1)
			(-1.6,0.15,#1)
			(-1.4,0.225,#1)
			(-1.2,0.3,#1)
			(-1.07,0.37,#1)
			(-1,0.5,#1)
			(-1,1,#1)
			(-1,1.25,#1)
			(-1,1.5,#1)
		};
	\draw[#2] (-2,0,#1) -- (-2,1.5,#1);
	\draw[#3] (-2,1.5,#1) -- (-1,1.5,#1);
	\draw[#3] (-2,1.125,#1) -- (-1,1.125,#1);
	\draw[#3] (-2,0.75,#1) -- (-1,0.75,#1);
	\draw[#3] (-2,0.375,#1) -- (-1.09,0.375,#1);
}

\begin{document}
\begin{tikzpicture}[
	x=\geomscale cm,
	y=\geomscale cm,
	z={(\zxcm cm,\zycm cm)},
	>={Stealth[length=3.2pt,width=2.4pt]},
	]

	% ----------------------------------------------------------
	% Far-field box: back plane and spanwise edges
	% ----------------------------------------------------------
	\draw[dotted,thick,black!45] (\aa,\hh,\zb) -- (\oo,\hh,\zb);
	\draw[dotted,thick,black!45] (\aa,0,\zb) -- (\aa,\hh,\zb);
	\draw[dotted,thick,black!45] (\oo,0,\zb) -- (\oo,\hh,\zb);
	\draw[dotted,thick,black!45] (\aa,\hh,0) -- (\aa,\hh,\zb);
	\draw[dotted,thick,black!45] (\oo,\hh,0) -- (\oo,\hh,\zb);

	% ----------------------------------------------------------
	% Gap and wall surfaces (painter's order: far to near)
	% ----------------------------------------------------------
	% gap floor
	\draw[fill=\floorcol,draw=\edgecol,thin]
	(0,-\dd,0) -- (\ww,-\dd,0) -- (\ww,-\dd,\zb) -- (0,-\dd,\zb) -- cycle;

	% upstream face inside the gap
	\draw[fill=\facecol,draw=\edgecol,thin]
	(0,0,0) -- (0,-\dd,0) -- (0,-\dd,\zb) -- (0,0,\zb) -- cycle;

	%downstream face of the gap
	\draw[fill=\facecol,draw=black, thin]
	(\ww,0,0) -- (\ww,-\dd,0) -- (\ww,-\dd,\zb) -- (\ww,0,\zb) -- cycle;

	% wall surface upstream of the gap
	\draw[fill=\surfcol,draw=\edgecol,thin]
	(\aa,0,0) -- (0,0,0) -- (0,0,\zb) -- (\aa,0,\zb) -- cycle;

	% wall surface downstream of the gap
	\draw[fill=\surfcol,draw=\edgecol,thin]
	(\ww,0,0) -- (\oo,0,0) -- (\oo,0,\zb) -- (\ww,0,\zb) -- cycle;

	% ----------------------------------------------------------
	% Inlet Blasius profile, swept over the whole span
	% ----------------------------------------------------------
	\begin{scope}[shift={(-4,0)}]

		% profile at z = b, and the lines ruling the swept surface
		\blprofile{\zb}{\ghostcol,thin}{-{Latex[length=4pt,width=3pt]},\ghostcol,thin}

		\draw[\ghostcol,thin] (-2,1.5,0) -- (-2,1.5,\zb);
		\draw[\ghostcol,thin] (-1,1.5,0) -- (-1,1.5,\zb);
		\draw[\ghostcol,thin] (-1,0.43,0) -- (-1,0.43,\zb);

		% profile at z = 0
		\blprofile{0}{black,thick}{-{Latex[length=5pt,width=4pt]},thick}

	\end{scope}

	% wall contour
	\draw[thin]
	(\aa,0) -- (0,0) -- (0,-\dd)
	-- (\ww,-\dd) -- (\ww,0) -- (\oo,0);

	% ----------------------------------------------------------
	% Displacement thickness of the smooth-wall reference flow at the
	% upstream edge of the gap: sketched at mid-span, not to scale
	% ----------------------------------------------------------
	\begin{scope}[shift={(0,0,\zb/2)}]
		\begin{scope}[yscale=1.7,xscale=0.8,shift={(1.85,0)}]

			\clip (-5,-0.2) rectangle (10,1);

			\def\col{orange!70!black}

			\draw[\col,thick,dotted,yscale=1.3]
			plot[smooth,tension=0.5] coordinates {
					(-2,0)
					(-1.8,0.075)
					(-1.6,0.15)
					(-1.4,0.225)
					(-1.2,0.3)
					(-1.07,0.37)
					(-1,0.5)
					(-1,1)
					(-1,1.25)
					(-1,1.5)
					(-1,1.75)
				};

			\draw[\col,thick,dotted] (-4,0) -- (-0.5,0);

			\draw[
			\col,
			{Latex[length=5pt,width=4pt]}-{Latex[length=5pt,width=4pt]},
			thick,
			dotted
			] (-2,0) -- (-2,0.75);

			\node[
				\col,
				font=\fsize,
				anchor=east
			] at (-2,0.325)
			{$\delta_{\mathrm{fp}}^*:=1$};

		\end{scope}
	\end{scope}

	% ----------------------------------------------------------
	% Origin of coordinates, at the upstream edge of the gap
	% ----------------------------------------------------------
	% filled: a white disc with a hairline outline was invisible against the
	% pale plate it sits on
	\draw[fill=white,thin] (0,0,0) circle (1.6pt);

	\node[font=\fsize,anchor=south east, xshift=0.1cm]
	at (0,0,0) {$\boldsymbol{x}=\boldsymbol{0}$};

	% ----------------------------------------------------------
	% Coordinate axes
	% ----------------------------------------------------------
	\begin{scope}[shift={(2,3.2)}]

		\draw[
			-Stealth,
			thick,
			green!60!black
		] (0,0.5) -- ++(0,1.4)
		node[
			above,
			anchor=east,
			font=\fsize
		] {$y$};

		\draw[
			-Stealth,
			thick,
			red!80!black
		] (0,0.5) -- ++(1.4,0)
		node[
			below,
			anchor=north west,
			font=\fsize
		] {$x$};

		\draw[
			-Stealth,
			thick,
			blue!70!black
		] (0,0.5) -- ++(0,0,2.5)
		node[
			anchor=north west,
			inner sep=1.5pt,
			font=\fsize
		] {$z$};

	\end{scope}

	% ----------------------------------------------------------
	% Dimensions
	% ----------------------------------------------------------
	% domain Omega
	\node[font=\fsizebig] (Omega)
	at (\aa*1.4,\hh,\zb*0.75) {$\Omega$};

	% width w
	\node[font=\fsize] (width)
	at (\ww/2,-\dd/2) {$w$};

	\draw[Stealth-]
	(0.1,-\dd/2) -- (width.west);

	\draw[-Stealth]
	(width.east) -- (\ww-0.1,-\dd/2);

	% depth d
	\draw[Stealth-Stealth]
	(-0.5,-0.1) -- (-0.5,-\dd);

	\node[font=\fsize,anchor=east]
	at (-0.6,-\dd/2) {$d$};

	% spanwise period b
	\node[font=\fsize] (span)
	at (\ww+1.9,0,\zb/2) {$b$};

	\draw[Stealth-]
	(\ww+1.9,0,0.12) -- (span.\zangle+180);

	\draw[-Stealth]
	(span.\zangle) -- (\ww+1.9,0,\zb-0.12);

	% height h
	\node[font=\fsize] (height)
	at (\aa-0.75,\hh/2) {$h$};

	\draw[-Stealth]
	(height.south) -- (\aa-0.75,0);

	\draw[Stealth-]
	(\aa-0.75,\hh) -- (height.north);

	% streamwise extents
	\def\deps{0.6}

	\node[font=\fsize] (ell1)
	at (\aa/2,-\dd-\deps) {$\ell_{\mathrm{i}}$};

	\draw[-Stealth]
	(ell1.east) -- (0,-\dd-\deps);

	\draw[Stealth-]
	(\aa,-\dd-\deps) -- (ell1.west);

	\node[font=\fsize] (ell2)
	at (\ww/2+\oo/2,-\dd-\deps) {$\ell_{\mathrm{o}}$};

	\draw[-Stealth]
	(ell2.east) -- (\oo,-\dd-\deps);

	\draw[Stealth-]
	(\ww,-\dd-\deps) -- (ell2.west);

	% ----------------------------------------------------------
	% Outer boundaries (near spanwise plane)
	% ----------------------------------------------------------
	\draw[dotted,thick] (\aa,\hh) -- (\oo,\hh);
	\draw[dotted,thick] (\aa,0) -- (\aa,\hh);
	\draw[dotted,thick] (\oo,0) -- (\oo,\hh);

	% ----------------------------------------------------------
	% Boundary conditions
	% ----------------------------------------------------------
	\node[font=\fsize] (W)
	at (9,-0.5) {Wall};

	\node[
	font=\fsize,
	anchor=south west,
	fill=white,
	inner sep=1.5pt
	] at (\aa,4.2)
	{$\boldsymbol{u}=\boldsymbol{u}_{\mathrm{BL}}(y)$};

	\node[
		font=\fsize,
		anchor=west
	] at (W.east)
	{$\boldsymbol{u}=\boldsymbol{0}$};

	\node[
		font=\fsize
	] at ({(\oo+\aa)/2},\hh,\zb/2)
	{$\boldsymbol{u}=\boldsymbol{u}_{\infty}(x)$};

	\node[
	font=\fsize,
	anchor=east,
	fill=white,
	inner sep=2pt,
	xshift=-2pt
	] at (\oo,\hh/2-1,\zb)
	{$
		\partial_{\boldsymbol{n}}\boldsymbol{u}
		=
		\boldsymbol{h}
		(\boldsymbol{u}_{\mathrm{out}})
	$};

	% the two spanwise boundaries are identified
	\def\xper{12}

	% witness lines lying in each of the two spanwise planes
	\draw[black,thin] (\xper,\hh-0.75,0) -- (\xper,\hh+0.15,0);
	\draw[black,thin] (\xper,\hh-0.75,\zb) -- (\xper,\hh+0.15,\zb);

	\draw[Stealth-Stealth] (\xper,\hh,0.1) -- (\xper,\hh,\zb-0.1);

	\node[
		font=\fsize,
		anchor=north west,
		inner sep=2.5pt
	] at (\xper,\hh,\zb/2) {periodic};

\end{tikzpicture}
\end{document}
